\documentclass{article}
\usepackage{amsmath}
\usepackage{amsfonts}
\usepackage{enumitem}
\usepackage{geometry}
\geometry{margin=1in}
\begin{document}

\section*{Homework 1 --- Solutions}

\subsection*{2. A noiseless 4-kHz channel is sampled every 1 msec. What is the maximum data rate?\\
How does the maximum data rate change if the channel is noisy, with a signal-to-noise ratio of 30 dB?}

\textbf{Answer:}

\textbf{Case 1: Noiseless Channel (Nyquist)}

Given:
\begin{itemize}[noitemsep]
  \item Bandwidth $B = 4000$ Hz
  \item Binary signal $\Rightarrow V = 2$ levels
\end{itemize}

Nyquist's theorem:
\[
\text{Max rate} = 2B \log_2 V
\]

Substituting:
\[
\text{Max rate} = 2 \cdot 4000 \cdot \log_2 2 = 8000 \text{ bps}
\]

\textbf{Case 2: Noisy Channel (Shannon)}

Shannon's theorem:
\[
\text{Max rate} = B \log_2\left(1 + \frac{S}{N} \right)
\]

Convert SNR from dB to power ratio:
\[
\frac{S}{N} = 10^{30/10} = 10^3 = 1000
\]

Now substitute into Shannon’s theorem:
\[
\text{Max rate} = 4000 \cdot \log_2(1 + 1000) = 4000 \cdot \log_2(1001)
\]

Approximating:
\[
\log_2(1001) \approx 9.97
\]

So:
\[
\text{Max rate} \approx 4000 \cdot 9.97 \approx 39880 \text{ bps}
\]

\subsection*{4. If a binary signal is sent over a 3-kHz channel whose signal-to-noise ratio is 20 dB, what is the maximum achievable data rate?}

\textbf{Answer:}

Given:
\begin{itemize}[noitemsep]
  \item Bandwidth $B = 3000$ Hz
  \item Binary signal $\Rightarrow V = 2$ levels
  \item SNR = 20 dB
\end{itemize}

Shannon’s theorem:
\[
\text{Max rate} = B \log_2\left(1 + \frac{S}{N} \right)
\]

Convert SNR from dB to power ratio:
\[
\frac{S}{N} = 10^{20/10} = 100
\]

Substitute:
\[
\text{Max rate} = 3000 \cdot \log_2(1 + 100) = 3000 \cdot \log_2(101)
\]

Approximating:
\[
\log_2(101) \approx 6.6582
\]

So:
\[
\text{Max rate} \approx 3000 \cdot 6.6582 = 19974.6 \text{ bps}
\]

\subsection*{6. What are the advantages of fiber optics over copper as a transmission medium? Is there any downside of using fiber optics over copper?}

\textbf{Answer:}

\textbf{Advantages:}
\begin{itemize}[noitemsep]
  \item Much higher bandwidth (THz range)
  \item Lower signal attenuation over long distances
  \item Immune to electromagnetic interference
  \item More secure and difficult to tap
\end{itemize}

\textbf{Disadvantages:}
\begin{itemize}[noitemsep]
  \item Higher installation cost
  \item More delicate; requires careful handling and alignment
\end{itemize}

\subsection*{13. Calculate the end-to-end transit time for a packet for both GEO (altitude: 35,800 km), MEO (altitude: 18,000 km), and LEO (altitude: 750 km) satellites.}

\textbf{Answer:}

We use the formula for propagation delay:
\[
 t = \frac{d}{c}
\]
Where:
\begin{itemize}
  \item $d$ = distance (one-way) in meters
  \item $c = 3 \times 10^8$ m/s (speed of light)
\end{itemize}

To compute round-trip time (RTT):
\[
\text{RTT} = 2 \cdot \frac{d}{c}
\]

\begin{itemize}
  \item \textbf{GEO:}
    \begin{itemize}
      \item $d = 35800$ km $= 35800 \cdot 1000 = 3.58 \times 10^7$ m
      \item $\text{RTT} = 2 \cdot \frac{3.58 \times 10^7}{3 \times 10^8} = 0.2387$ s $= \boxed{239 \text{ ms}}$
    \end{itemize}

  \item \textbf{MEO:}
    \begin{itemize}
      \item $d = 18000$ km $= 18000 \cdot 1000 = 1.8 \times 10^7$ m
      \item $\text{RTT} = 2 \cdot \frac{1.8 \times 10^7}{3 \times 10^8} = 0.12$ s $= \boxed{120 \text{ ms}}$
    \end{itemize}

  \item \textbf{LEO:}
    \begin{itemize}
      \item $d = 750$ km $= 750 \cdot 1000 = 7.5 \times 10^5$ m
      \item $\text{RTT} = 2 \cdot \frac{7.5 \times 10^5}{3 \times 10^8} = 0.005$ s $= \boxed{5 \text{ ms}}$
    \end{itemize}
\end{itemize}

\subsection*{16. Prove that in 4B/5B encoding, a signal transition will occur at least every four bit times.}

\textbf{Answer:}

In 4B/5B encoding, each 4-bit group is mapped to a 5-bit codeword with no more than three consecutive 0s.
This ensures that at least one signal transition occurs every 5-bit transmission, i.e., every 4 original data bits, improving clock synchronization.

\subsection*{24. How many frequencies does a full-duplex QAM-64 modem use?}

\textbf{Answer:}

A full-duplex modem uses separate frequency bands for transmitting and receiving simultaneously. Thus, a full-duplex QAM-64 modem uses:
\[
\boxed{2 \text{ frequencies}}
\]

\subsection*{30. What is the difference, if any, between the demodulator part of a modem and the coder part of a codec? (After all, both convert analog signals to digital ones.)}

\textbf{Answer:}

\begin{itemize}
  \item A \textbf{demodulator} extracts digital data from a modulated analog carrier signal. It operates in data communication systems (modems).
  \item A \textbf{coder} in a codec converts analog voice or audio signals into a digital stream for storage or transmission.
\end{itemize}

While both perform analog-to-digital conversion, they serve different applications: data vs. voice/audio, and differ in method and context.

\subsection*{37. Three packet-switching networks each contain $n$ nodes. The first network has a star topology with a central switch, the second is a (bidirectional) ring, and the third is fully interconnected, with a wire from every node to every other node. What are the best-, average-, and worst-case transmission paths in hops?}

\textbf{Answer:}

\begin{itemize}
  \item \textbf{Star:} Best = 1, Average = 2, Worst = 2
  \item \textbf{Ring:} Best = 1, Average = $n/4$, Worst = $n/2$
  \item \textbf{Fully-connected:} Best = Average = Worst = 1
\end{itemize}

\subsection*{44. Suppose that A, B, and C are simultaneously transmitting 0 bits, using a CDMA system with the chip sequences of Fig. 2-28(a). What is the resulting chip sequence?}

\textbf{Answer:}

Each 0 bit is encoded as the inverse of the chip sequence.

\begin{itemize}
  \item $A = (+1, -1, +1, -1) \Rightarrow$ sends $(-1, +1, -1, +1)$
  \item $B = (+1, +1, -1, -1) \Rightarrow$ sends $(-1, -1, +1, +1)$
  \item $C = (-1, +1, +1, -1) \Rightarrow$ sends $(+1, -1, -1, +1)$
\end{itemize}

Add corresponding chips:
\[
(-1)+(-1)+(+1) = -1,\quad (+1)+(-1)+(-1) = -1,\quad (-1)+(+1)+(-1) = -1,\quad (+1)+(+1)+(+1) = +3
\]

\[
\text{Resulting chip sequence: } \boxed{(-1, -1, -1, +3)}
\]

\subsection*{46. A CDMA receiver gets the following chips: $(-1, +1, -3, +1, -1, -3, +1, +1)$.
 Assuming the chip sequences defined in Fig. 2-28(a), which stations transmitted, and which bits did each one send?}

\textbf{Answer:}

Given chip sequence is 8 chips: divide into two 4-chip segments.

\[
\text{C1} = (-1, +1, -3, +1),\quad \text{C2} = (-1, -3, +1, +1)
\]

For each user (A, B, C), compute dot product with C1 and C2. If result is negative → 0; if positive → 1.

\textbf{User A:} $(+1, -1, +1, -1)$
\[
\text{C1} \cdot A = (-1)(+1) + (+1)(-1) + (-3)(+1) + (+1)(-1) = -1 -1 -3 -1 = -6 \Rightarrow \text{bit 0}
\]
\[
\text{C2} \cdot A = (-1)(+1) + (-3)(-1) + (+1)(+1) + (+1)(-1) = -1 +3 +1 -1 = 2 \Rightarrow \text{bit 1}
\]

\textbf{User B:} $(+1, +1, -1, -1)$
\[
\text{C1} \cdot B = -1 +1 +3 -1 = 2 \Rightarrow \text{bit 1}
\]
\[
\text{C2} \cdot B = -1 -3 -1 -1 = -6 \Rightarrow \text{bit 0}
\]

\textbf{User C:} $(-1, +1, +1, -1)$
\[
\text{C1} \cdot C = +1 +1 -3 -1 = -2 \Rightarrow \text{bit 0}
\]
\[
\text{C2} \cdot C = +1 -3 +1 -1 = -2 \Rightarrow \text{bit 0}
\]

\textbf{Result:}
\begin{itemize}
  \item A sent: 0, 1
  \item B sent: 1, 0
  \item C sent: 0, 0
\end{itemize}

\end{document}
